\section{Objectifs}

\begin{itemize}
\item Placer et lire un réel sur une droite graduée.
\item Distinguer nombres décimaux, rationnels, irrationnels et réels.
\item Représenter un intervalle et tester une appartenance.
\item Utiliser la valeur absolue uniquement comme distance dans le cadre du programme de seconde.
\item Résoudre et interpréter une inéquation de la forme \texttt{|x-a|≤r}.
\item Donner un encadrement décimal d'amplitude imposée et arrondir avec un nombre de chiffres significatifs adapté à la situation.
\end{itemize}

\section{Prérequis}

\begin{itemize}
\item droite graduée
\item nombres relatifs
\item comparaison de nombres
\item puissances de 10
\end{itemize}

\section{Activité de découverte}

Une application de géolocalisation donne une position à \texttt{0,1 km} près. On cherche tous les emplacements possibles sur une ligne droite : cela conduit à un intervalle centré et à une écriture avec une distance.

\section{Vocabulaire}

\begin{itemize}
\item \textbf{réel} \texttt{ℝ}, \textbf{décimal} \texttt{𝔻}, \textbf{rationnel} \texttt{ℚ}, \textbf{irrationnel}
\item \textbf{droite numérique}, \textbf{intervalle}, \textbf{borne}, \textbf{amplitude}
\item \textbf{valeur absolue}, \textbf{distance}, \textbf{encadrement}, \textbf{arrondi}, \textbf{chiffres significatifs}
\end{itemize}

\section{Cours}

\subsection{Les ensembles de nombres}

\begin{itemize}
\item \texttt{𝔻} désigne les nombres décimaux, c'est-à-dire ceux qui possèdent une écriture décimale finie.
\item \texttt{ℚ} désigne les nombres rationnels, qui peuvent s'écrire comme quotient de deux entiers avec dénominateur non nul.
\item \texttt{ℝ} désigne les nombres réels, représentés par les points de la droite numérique.
\item Des nombres comme \texttt{√2} et \texttt{π} sont irrationnels : ils appartiennent à \texttt{ℝ} mais pas à \texttt{ℚ}.
\end{itemize}

Une calculatrice ne stocke qu'une approximation d'un réel général : il faut distinguer le nombre mathématique et sa valeur affichée.

\subsection{Intervalles}

Un intervalle est un ensemble de réels défini par des inégalités. Les crochets indiquent si une borne finie appartient ou non à l'intervalle.

Exemples :

\begin{itemize}
\item \texttt{[a;b]} correspond à \texttt{a≤x≤b} ;
\item \texttt{]a;b[} correspond à \texttt{a<x<b} ;
\item \texttt{[a;+∞[} correspond à \texttt{x≥a}.
\end{itemize}

\subsection{Valeur absolue comme distance}

En seconde, la valeur absolue est introduite pour exprimer une distance :

\begin{itemize}
\item \texttt{|a|} est la distance de \texttt{a} à \texttt{0} ;
\item \texttt{|a-b|} est la distance entre \texttt{a} et \texttt{b}.
\end{itemize}

Pour \texttt{r≥0} :

\texttt{|x-a|≤r ⇔ x∈[a-r;a+r]}.

Cette équivalence permet de passer d'une contrainte de distance à un intervalle et inversement.

\subsection{Encadrements décimaux et précision}

Un encadrement décimal à \texttt{10^-n} près doit avoir une amplitude inférieure ou égale à la précision demandée. Dans un problème concret, l'arrondi doit utiliser un nombre de chiffres significatifs cohérent avec les données et la grandeur étudiée.

Exemple : \texttt{1,41<√2<1,42} est un encadrement d'amplitude \texttt{0,01}.

\section{Exemples corrigés}

\begin{itemize}
\item \texttt{|x-3|≤2} signifie que la distance de \texttt{x} à \texttt{3} est au plus \texttt{2}, donc \texttt{x∈[1;5]}.
\item \texttt{1/3} est rationnel mais n'est pas décimal.
\item \texttt{√2} est réel et irrationnel.
\item Une longueur calculée à partir de mesures données au millimètre ne doit pas être affichée avec une précision artificielle de dix décimales.
\end{itemize}

\coursefigure{/images/mathematiques/latex-migration/2nde-nombres-reels-intervalles-figure-1.svg}{Droite graduée montrant les solutions de valeur absolue entre 1 et 5.}{Une contrainte de distance peut se traduire par un intervalle centré.}

\section{Algorithmique en Python}

\subsection{Encadrer \texttt{√2} par balayage}

\begin{verbatim}
def encadre_racine2(pas):
    x = 1
    while x * x < 2:
        x = x + pas
    return x - pas, x

print(encadre_racine2(0.01))
\end{verbatim}

Le programme avance par pas réguliers jusqu'à encadrer \texttt{√2}. Le pas choisi contrôle l'amplitude de l'encadrement.

\section{Démonstrations à travailler}

\begin{itemize}
\item \texttt{1/3} est rationnel mais n'est pas un nombre décimal.
\item \texttt{√2} est irrationnel.
\end{itemize}

\section{Erreurs fréquentes}

\begin{itemize}
\item Confondre une borne ouverte et une borne fermée.
\item Utiliser la valeur absolue pour d'autres usages non attendus en seconde au lieu de l'interpréter comme distance.
\item Écrire un encadrement sans vérifier son amplitude.
\item Confondre nombre réel et approximation affichée par la calculatrice.
\item Donner trop de chiffres après la virgule sans lien avec la précision des données.
\end{itemize}

\section{Formules et idées à retenir}

\begin{itemize}
\item \texttt{|a-b|} est la distance entre \texttt{a} et \texttt{b}.
\item \texttt{|x-a|≤r ⇔ x∈[a-r;a+r]} pour \texttt{r≥0}.
\item \texttt{𝔻⊂ℚ⊂ℝ}.
\end{itemize}

\section{Synthèse}

Les nombres réels sont étudiés à travers la droite numérique, les intervalles et l'approximation. La valeur absolue garde ici un sens géométrique de distance ; encadrer et arrondir exige de contrôler explicitement la précision du résultat.
